\subsubsection*{HyperKvasir}

The HyperKvasir dataset is a comprehensive, multi-purpose gastrointestinal dataset designed to support a wide range of computer vision tasks, including classification, segmentation, and video analysis. It consists of 110{,}079 images, of which 10{,}662 are labeled across 23 classes relevant to gastroenterology. In addition, it includes 1{,}000 manually segmented polyp images and 374 labeled videos comprising 889{,}372 frames.

The dataset captures a wide diversity of gastrointestinal findings—ranging from polyps and esophagitis to bleeding and mucosal inflammation—collected from real clinical procedures at a Norwegian hospital using Olympus and Pentax endoscopes. Labels were initially assigned by junior medical staff and subsequently validated by experienced gastroenterologists, ensuring medical accuracy.

The image classification baseline was established using deep convolutional neural networks including ResNet-50, ResNet-152, and DenseNet-161. Evaluation metrics included macro- and micro-averaged F1-scores and the Matthews Correlation Coefficient (MCC), with the best-performing ensemble reaching an MCC of 0.902 and micro-F1 of 0.910.

A key strength of HyperKvasir lies in its broad coverage of anatomical and pathological categories and its modular structure, allowing flexible use for supervised and semi-supervised learning. It has become a standard benchmark in gastrointestinal AI research and supports the development of clinically relevant decision support systems beyond polyp detection.
